\documentclass[aspectratio=169]{beamer}
\usepackage[no-math,deluxe,haranoaji]{luatexja-preset}
\renewcommand{\kanjifamilydefault}{\gtdefault}
\renewcommand{\emph}[1]{{\upshape\bfseries #1}}
\usetheme{metropolis}
\metroset{block=fill}
\setbeamertemplate{navigation symbols}{}
\setbeamertemplate{blocks}[rounded][shadow=false]
\usecolortheme[rgb={0.7,0.2,0.2}]{structure}
%%%%%%%%%%%%%%%%%%%%%%%%%%%
\usepackage{media9}
%%%%%%%%%%%%%%%%%%%%%%%%%%%
%% さまざまなアイコン
%%%%%%%%%%%%%%%%%%%%%%%%%%%
\usepackage{fontawesome}
%\usepackage{figchild}
\usepackage{twemojis}
\usepackage{utfsym}
\usepackage{bclogo}
\usepackage{marvosym}
\usepackage{fontmfizz}
%%%%%%%%%%%%%%%%%%%%%%%%%%%
\usepackage{tikz}
\usetikzlibrary{backgrounds}
\usepackage{tcolorbox}
\usepackage{tikzpeople}
\usepackage{tikzducks}
\usepackage{xcolor}
\usepackage{amsmath}
\usepackage{tipa}
\usepackage{circledsteps}
%%%%%%%%%%%%%%%%%%%%%%%%%%%
%% 場合分け
\usepackage{cases}
%%%%%%%%%%%%%%%%%%%%%%%%%%%
% \myAnch{<名前>}{<色>}{<テキスト>}
% 指定のテキストを指定の色の四角枠で囲み, 指定の名前をもつTikZの
% ノードとして出力する. 図には remeber picture 属性を付けている
% ので外部から参照可能である.
\newcommand*{\myAnch}[3]{%
  \tikz[remember picture,baseline=(#1.base)]
    \node[draw,rectangle,#2] (#1) {\normalcolor #3};
}
%%%%%%%%%%%%%%%%%%%%%%%%%%%%
%% 音声リンク表示
\newcommand{\myaudio}[1]{\href{#1}{\faVolumeUp}}
%%%%%%%%%%%%%%%%%%%%%%%%%%%
% \myEmph コマンドの定義
%\newcommand{\myEmph}[3]{%
%    \textbf<#1>{\color<#1>{#2}{#3}}%
%}
\usepackage{xparse} % xparseパッケージの読み込み
\NewDocumentCommand{\myEmph}{O{} m m}{%
    \def\argOne{#1}%
    \ifx\argOne\empty
        \textbf{\color{#2}{#3}}% オプション引数が省略された場合
    \else
        \textbf<#1>{\color<#1>{#2}{#3}}% オプション引数が指定された場合
    \fi
}
%%%%%%%%%%%%%%%%%%%%%%%%%%%
%% 文末の上昇イントネーション記号 \myRisingPitch
%% 通常のイントネーション \myDownwardPitch
%% https://note.com/dan_oyama/n/n8be58e8797b2
%%%%%%%%%%%%%%%%%%%%%%%%%%%
\newcommand{\myRisingPitch}{
\begin{tikzpicture}[scale=0.3,baseline=0.3]
\draw[->,>=stealth] (0,0) to[bend right=45] (1,1);
\end{tikzpicture}
}
\newcommand{\myDownwardPitch}{
\begin{tikzpicture}[scale=0.3,baseline=0.3]
\draw[->,>=stealth] (0,1) to[bend left=45] (1,0);
\end{tikzpicture}
}
%%%%%%%%%%%%%%%%%%%%%%%%%%%
\title{English is fun.}
\subtitle{Are you hungry?}
\author{}
\institute[]{}
\date[]

%%%%%%%%%%%%%%%%%%%%%%%%%%%%
%% TEXT
%%%%%%%%%%%%%%%%%%%%%%%%%%%%
\begin{document}
%%%%%%%%%%%%%%%%%%%%%%%%%%%%
\begin{frame}[plain]
  \titlepage
\end{frame}
%%%%%%%%%%%%%%%%%%%%%%%%%%%%
\section*{授業の流れ}
%%%%%%%%%%%%%%%%%%%%%%%%%%%%
\begin{frame}[plain]
  \frametitle{授業の流れ}
  \tableofcontents
\end{frame}
%%%%%%%%%%%%%%%%%%%%%%%%%%%%
\section{be動詞の疑問文}
\subsection{一人称}
%%%%%%%%%%%%%%%%%%%%%%%%%%%%
\begin{frame}[plain,t]{一人称}
 \large

疑問文のつくり方($\text{主語} = \text{\alt<1-5>{I}{I $\rightarrow$ We}}$)%
\hfill{\tiny 0102}\,{\scriptsize\myaudio{./audio/008_question_be_01.mp3}}

\vspace{10pt}

\myAnch{s-1}{orange}{\alt<1-5>{I}{We}} \myAnch{be-1}{olive}{\alt<1-5>{am}{are}} late for the movie.%
\hfill{\scriptsize be late for ～に遅れる}

\vspace{25pt}

\visible<2->{\myAnch{be-2}{olive}{\alt<1-5>{Am}{Are}} \myAnch{s-2}{orange}{\alt<1-5>{I}{we}} late for the movie\myAnch{question}{orange}{?}}\visible<5->{\scalebox{2}{\myRisingPitch}}

\visible<4->{{\small be動詞を先頭に}}\hspace{70pt}\visible<5->{{\small ?と最後のイントネーションに注意}}


\visible<3->{\begin{tikzpicture}[remember picture, overlay]
 \draw[line width=2pt, opacity=.5, orange, ->] (s-1.south) to[out=-90, in=90] (s-2.north);
 \draw[line width=2pt, opacity=.5, olive, ->] (be-1.south) to[out=-90, in=90] (be-2.north);
\end{tikzpicture}}

\begin{block}<7->{Topics for Today}\small
 \begin{itemize}\setbeamertemplate{items}[square]
  \item \textbf{I am} ～ . の疑問文は\textbf{Am I} ～ ?
  \item \textbf{We are} ～ . の疑問文は\textbf{Are we} ～ ?%
	\hfill{\scriptsize \textbf{Are we} late for the movie?}
  \item \textbf{be}動詞を先頭にもっていくだけです
\end{itemize}
\end{block}
\end{frame}
%%%%%%%%%%%%%%%%%%%%%%%%%%%%%%%
\subsection{二人称}
%%%%%%%%%%%%%%%%%%%%%%%%%%%%%%%
\begin{frame}[plain,t]{二人称}
 \large

疑問文のつくり方($\text{主語} = \text{You}$)%
\hfill{\tiny 0057}\,{\scriptsize \myaudio{./audio/008_question_be_02.mp3}}


\vspace{10pt}

\myAnch{s-1}{orange}{You} \myAnch{be-1}{olive}{are} happy.\scalebox{2}{\myDownwardPitch}
\vspace{20pt}

\visible<2->{\myAnch{be-2}{olive}{Are} \myAnch{s-2}{orange}{you} happy\myAnch{question}{orange}{?}}\visible<5->{\scalebox{2}{\myRisingPitch}}

\visible<4->{\small be動詞を先頭に}\hspace{52pt}\visible<5->{\small ?と最後のイントネーションに注意}

\visible<3->{\begin{tikzpicture}[remember picture, overlay]
%\draw [help lines] (0,0) grid (10,4);%(0,0)から(10,4)までの"細線の方眼"
 \draw[line width=2pt, opacity=.5, orange, ->] (s-1.south) to[out=-90, in=90] (s-2.north);
 \draw[line width=2pt, opacity=.5, olive, ->] (be-1.south) to[out=-90, in=90] (be-2.north);
\end{tikzpicture}}

\begin{block}<6->{Topics for Today}\small
 \begin{itemize}\setbeamertemplate{items}[square]
  \item \textbf{You are} ～ . の疑問文は\textbf{Are you} ～ ?
  \item \textbf{be}動詞を先頭にもっていくだけです
\end{itemize}
\end{block}
\end{frame}
%%%%%%%%%%%%%%%%%%%%%%%%%%%%%%%
\subsection{三人称単数}
%%%%%%%%%%%%%%%%%%%%%%%%%%%%%%%
\begin{frame}[t,plain]{三人称単数}
 \large

疑問文のつくり方($\text{主語} = \text{\temporal<6>{She}{She $\rightarrow$ He}{She $\rightarrow$ He $\rightarrow$ Our teacher}}$)%
\hfill{\tiny 0225}\,{\scriptsize \myaudio{./audio/008_question_be_03.mp3}}

\myAnch{s-1}{orange}{\temporal<6>{She}{He}{Our teacher}} \myAnch{be-1}{olive}{is} from Australia.\hfill{\scriptsize be from ～の出身だ}\\
\hfill{\scriptsize Australia \textipa{/O:str\'eIlj@/} オーストラリア}

\vspace{20pt}

\visible<2->{\myAnch{be-2}{olive}{Is} \myAnch{s-2}{orange}{\temporal<6>{she}{he}{our teacher}} from Australia\myAnch{question}{orange}{?}}\visible<5->{\scalebox{2}{\myRisingPitch}}

\visible<4->{{\small be動詞を先頭に}}\hspace{90pt}\visible<5->{{\small ?と最後のイントネーションに注意}}

\visible<3->{\begin{tikzpicture}[remember picture, overlay]
 \draw[line width=2pt, opacity=.5, orange, ->] (s-1.south) to[out=-90, in=90] (s-2.north);
 \draw[line width=2pt, opacity=.5, olive, ->] (be-1.south) to[out=-90, in=90] (be-2.north);
\end{tikzpicture}}

\vspace{-5pt}

\begin{block}<8->{Topics for Today}\small
 \begin{itemize}\setbeamertemplate{items}[square]
  \item \textbf{She(He, Our teacher) is} ～ . の疑問文は\textbf{Is she(he, our teacher)} ～ ?
  \item \textbf{be}動詞を先頭にもっていくだけです
\end{itemize}
\end{block}

\vspace{-5pt}

\hfill\visible<7->{{\scriptsize teacher $<$ teach $+$ er($=$ ～する人)}}

\end{frame}
%%%%%%%%%%%%%%%%%%%%%%%%%%%%%%%%%%%%%
\subsection{三人称複数}
%%%%%%%%%%%%%%%%%%%%%%%%%%%%%%%%%%%%%
\begin{frame}[plain,t]{三人称複数}
 \large

疑問文のつくり方($\text{主語} = \text{\alt<1-5>{They}{They $\rightarrow$ Jack and Betty}}$)%
\hfill{\tiny 0146}\,{\scriptsize\myaudio{./audio/008_question_be_04.mp3}}

\vspace{10pt}

\myAnch{s-1}{orange}{\alt<1-5>{They}{Jack and Betty}} \myAnch{be-1}{olive}{are} good singers.%
\hfill{\scriptsize singer \textipa{/s\'IN\textrhookschwa/} 歌手 $<$ sing $+$ er}\\
\hfill{\scriptsize er ～する人}

\vspace{25pt}

\visible<2->{\myAnch{be-2}{olive}{Are} \myAnch{s-2}{orange}{\alt<1-5>{they}{Jack and Betty}} good singers\myAnch{question}{orange}{?}}\visible<5->{\scalebox{2}{\myRisingPitch}}

\visible<4->{{\small be動詞を先頭に}}\hspace{70pt}\visible<5->{{\small ?と最後のイントネーションに注意}}


\visible<3->{\begin{tikzpicture}[remember picture, overlay]
 \draw[line width=2pt, opacity=.5, orange, ->] (s-1.south) to[out=-90, in=90] (s-2.north);
 \draw[line width=2pt, opacity=.5, olive, ->] (be-1.south) to[out=-90, in=90] (be-2.north);
\end{tikzpicture}}

\begin{block}<7->{Topics for Today}\small
 \begin{itemize}\setbeamertemplate{items}[square]
  \item \textbf{They(Jack and Betty) are} ～ . の疑問文は\textbf{Are they(Jack and Betty)} ～ ?
  \item \textbf{be}動詞を先頭にもっていくだけです
\end{itemize}
\end{block}
\end{frame}
%%%%%%%%%%%%%%%%%%%%%%%%%%%%%%%%
%%%%%%%%%%%%%%%%%%%%%%%%%%%%%%%%%%%%%
\section{確認問題}
%%%%%%%%%%%%%%%%%%%%%%%%%%%%%%%%%%%%%
\begin{frame}[plain]\frametitle{Exercises}
つぎの各文を疑問文にしましょう%
\hfill{\tiny 0340}\,{\scriptsize\myaudio{./audio/008_question_be_05.mp3}}

 \begin{enumerate}
  \item I am late for the meeting.
        \onslide<5->{$\longrightarrow$ Am I late for the meeting?}%
	\hfill{\scriptsize be late for ～に遅れる}
  \item You are hungry.
        \onslide<6->{$\longrightarrow$ Are you hungry?}%
	\hfill{\scriptsize hungry \textipa{/h\'2Ngri/} 空腹の}
  \item They are happy.
        \onslide<7->{$\longrightarrow$ Are they happy?}
  \item She is in the kitchen.
        \onslide<8->{$\longrightarrow$ Is she in the kitchen?}%
	\hfill{\scriptsize kitchen \textipa{/k\'ItS@n/} 台所、キッチン}
  \item Our teacher is busy.
        \onslide<9->{$\longrightarrow$ Is our teacher busy?}%
	\hfill{\scriptsize busy \textipa{/b\'Izi/} 忙しい}
 \end{enumerate}

\begin{block}<2->{Topics for Today}\small
be動詞の疑問文のつくり方
\begin{itemize}\setbeamertemplate{items}[square]
 \item<3->\ \textbf{be}動詞を先頭にする（主語はその後ろ）\\
\hfill\Circled[fill color = white]{\,\,$\text{S}+\text{be動詞}\,\,$\,\,\ldots\,\,\,}\,\,.\,\,$\longrightarrow$\,\,\,\,\Circled[fill color = white]{\,\,$\text{be動詞}+\text{S}\,\,$\,\,\ldots\,\,\,}\,\,? 
 \item<4-> 文末に`?'をつける　　　イントネーションは\myRisingPitch
\end{itemize}
      \end{block}
\end{frame}
%%%%%%%%%%%%%%%%%%%%%%%%%%%%%%%
\begin{frame}[plain]\frametitle{Exercises}

{\small 日本文の意味になるよう語を並べ替えましょう。先頭に来る語は大文字ではじめてください。クエスチョンマークは補ってください}%
\hfill{\tiny 0321}\,{\scriptsize\myaudio{./audio/008_question_be_06.mp3}}

\begin{enumerate}
  \item あなたは幸せですか。 ( happy / you / are )\hfill\visible<2->{Are you happy?}
  \item 彼女は学生ですか。( a / she / student / is )\hfill\visible<3->{Is she a student?}
  \item 彼らは日本の出身ですか。( from / are / they / Japan )\hfill\visible<4->{Are they from Japan?}
  \item 彼はあなたのお父さんですか。( father / he / your / is )\hfill\visible<5->{Is he your father?}
  \item 私たちは会議に遅刻でしょうか。( for / are / late / we / the meeting )\\
\hfill\visible<6->{Are we late for the meeting?}
  \item きょうは寒いですか。( cold / today / it / is )\hfill\visible<7->{Is it cold today?}
  \item その犬たちは空腹ですか。( dogs / hungry / the / are )\hfill\visible<8->{Are the dogs hungry?}
  \item トムは部屋にいますか。( room / is / in / Tom / the )
\hfill\visible<9->{Is Tom in the room?}
\end{enumerate}
\end{frame}
%%%%%%%%%%%%%%%%%%%%%%%%%%%%%%%
\section{be動詞の疑問文（まとめ）}
%%%%%%%%%%%%%%%%%%%%%%%%%%%%%%%
\begin{frame}[plain]{要点}

\begin{block}{be動詞を用いた文の疑問文のつくり方}

\begin{itemize}\setbeamertemplate{items}[square]
 \item   \textbf{be}動詞を先頭にする
 \item   文末に`?'をつける\hfill{\scriptsize `?'をquestion mark（クエスチョンマーク）といいます}

 \item   イントネーションは\myRisingPitch
\end{itemize}
     \end{block}
\hfill{\scriptsize 主語がなんであろうが、疑問文のつくり方は同じです}%
{\begin{tikzpicture}
	\duck[strawhat=brown!50!white,
ribbon=gray];
       \end{tikzpicture}}

\hfill{\tiny audio\_overview 1755}\,{\scriptsize \myaudio{./audio/overview/008_question_be_audio_overview.mp4}}

\end{frame}
%%%%%%%%%%%%%%%%%%%%%%%%%%%%%%%%
\end{document}

